
\sgaStableTarget{sga6.e09.sec.3}\subsection*{3. Projective bundles and flag bundles}

\sgaStableTarget{sga6.e09.prop.3-1}\textbf{Proposition 3.1.} \uline{Let} @@M0@@ \uline{be a noetherian scheme}, @@M1@@ \uline{a locally free} @@M2@@\uline{-Module of rank} @@M3@@, @@M4@@ \uline{a sequence of integers such that} @@M5@@, \uline{and} @@M6@@ \uline{the scheme of flags of type} @@M7@@ \uline{in} @@M8@@. \uline{Then the canonical homomorphism}:
@@M9@@
\uline{is an isomorphism}.

Resume the notation of VI 4.5. Thus @@M10@@ denotes the canonical quotients of the inverse image of @@M11@@ on @@M12@@, @@M13@@ the kernel of @@M14@@ (@@M15@@), and @@M16@@ the class in @@M17@@ of @@M18@@. Then the @@M19@@ generate @@M20@@ as a @@M21@@-algebra (VI 4.6). Likewise, the images of the @@M22@@ in @@M23@@ generate @@M24@@ as an algebra over @@M25@@ for every point @@M26@@ of @@M27@@. Since @@M28@@ is a flag scheme over a field, it is a regular scheme, so the homomorphism @@M29@@ is an isomorphism, and thus the homomorphism @@M30@@ is surjective. By 1.3, this implies that the homomorphism:
@@M31@@
is surjective.

To prove injectivity, one may first reduce to the case of the complete flag bundle @@M32@@. Indeed, suppose that
@@M33@@
is injective. Since @@M34@@ factors through @@M35@@, this

